% Generated by scripts/render_editions.py from data/problem-statements.json.
\expandafter\def\csname knauthor@4.55\endcsname{K. W. Roggenkamp}
\expandafter\def\csname knquestion@4.55\endcsname{For a finite group $G$ and the localization $\Z_{(p)}$ at a prime $p$, does every projective $\Z_{(p)}G$-module have a decomposition into indecomposable projective modules unique up to permutation and isomorphism?}
\expandafter\def\csname knscope@4.55\endcsname{}
\expandafter\def\csname knauthor@4.75\endcsname{V. P. Shunkov}
\expandafter\def\csname knquestion@4.75\endcsname{Let $G$ be periodic, with an involution $i$, and with Sylow $2$-subgroups locally cyclic or generalized quaternion. Must $iO_{2\prime}(G)$ be central in $G/O_{2\prime}(G)$?}
\expandafter\def\csname knscope@4.75\endcsname{}
\expandafter\def\csname knauthor@6.47\endcsname{C. D. H. Cooper; communicated by E. I. Khukhro}
\expandafter\def\csname knquestion@6.47\endcsname{If a binary group word $v$ defines a new group operation $x\odot y=v(x,y)$ on a group $G$, must the resulting group $G_v$ belong to the variety generated by $G$?}
\expandafter\def\csname knscope@6.47\endcsname{\noindent\emph{Scope.} The proof below covers nilpotency class at most two.}
\expandafter\def\csname knauthor@9.4\endcsname{A. A. Gvaramiya}
\expandafter\def\csname knquestion@9.4\endcsname{If a variety, quasivariety or pseudovariety $\mathcal M$ of groups is generated by one finite group, must the class $I\mathcal M$ of quasigroups isotopic to groups in $\mathcal M$ be generated by one finite quasigroup?}
\expandafter\def\csname knscope@9.4\endcsname{\noindent\emph{Scope.} The argument below concerns the ordinary one-sorted variety of quasigroups, with its stated signature and closure operations.}
\expandafter\def\csname knauthor@9.45\endcsname{Yu. D. Popov, I. V. Protasov}
\expandafter\def\csname knquestion@9.45\endcsname{For $a\in\Q^n$, find necessary and sufficient conditions on its coordinates for $S(a)=\Z a+\Z^n\subseteq\R^n$ to have a basis orthogonal for the standard scalar product.}
\expandafter\def\csname knscope@9.45\endcsname{}
\expandafter\def\csname knauthor@9.47\endcsname{I. V. Protasov}
\expandafter\def\csname knquestion@9.47\endcsname{Does every scattered compact space embed homeomorphically in the space, with its $E$-topology, of all closed noncompact subgroups of a suitable locally compact group?}
\expandafter\def\csname knscope@9.47\endcsname{}
\expandafter\def\csname knauthor@10.32\endcsname{Yu. I. Merzlyakov}
\expandafter\def\csname knquestion@10.32\endcsname{For nonzero integers $r,s$ with at least one odd, find a logarithmic bound in the prime-divisor parameter $m(r,s)$ ensuring that $x^ry^s$ takes every value in $S_n$. The analogous alternating-group bound has the form $n\geq\max\{N_0(c),c\log m(r,s)\}$.}
\expandafter\def\csname knscope@10.32\endcsname{\noindent\emph{Scope.} The convention for unit exponents and the parameter are specified below.}
\expandafter\def\csname knauthor@10.35\endcsname{G. A. Noskov}
\expandafter\def\csname knquestion@10.35\endcsname{Is every finitely generated torsion-free subgroup of $\GL_n(\C)$ residually in the class of torsion-free subgroups of $\GL_n(\overline{\Q})$, with the same $n$?}
\expandafter\def\csname knscope@10.35\endcsname{\noindent\emph{Scope.} The claimed answer is withdrawn; the following rational-field observation does not answer this question.}
\expandafter\def\csname knauthor@10.62\endcsname{A. I. Sozutov}
\expandafter\def\csname knquestion@10.62\endcsname{Construct a periodic group without a subgroup of index two, generated by a conjugacy class $X$ of involutions, such that the product of any two members of $X$ has odd order.}
\expandafter\def\csname knscope@10.62\endcsname{}
\expandafter\def\csname knauthor@11.115\endcsname{V. P. Shaptala}
\expandafter\def\csname knquestion@11.115\endcsname{Let $X$ be a noncyclic free group, $1\ne N\trianglelefteq X$, and $N\leq T<X$. If $N$ is not maximal in $X$, must $[T,N]$ be a proper subgroup of $[X,N]$?}
\expandafter\def\csname knscope@11.115\endcsname{}
\expandafter\def\csname knauthor@11.116\endcsname{L. N. Shevrin}
\expandafter\def\csname knquestion@11.116\endcsname{Does every Chernikov group containing no direct product of two quasicyclic groups for the same prime have a subgroup lattice of finite order dimension? The order dimension is the least number of linear orders whose intersection is the given partial order.}
\expandafter\def\csname knscope@11.116\endcsname{}
\expandafter\def\csname knauthor@12.40\endcsname{A. S. Kondratiev}
\expandafter\def\csname knquestion@12.40\endcsname{For an irreducible $p$-modular character $\varphi$ of a finite group $G$, find the best bound $\varphi(1)_p\leq f(|G|_p)$, where the subscript denotes the $p$-part.}
\expandafter\def\csname knscope@12.40\endcsname{}
\expandafter\def\csname knauthor@12.69\endcsname{K. N. Ponomar\"ev}
\expandafter\def\csname knquestion@12.69\endcsname{Let $F$ be a countably infinite field, $G$ a finite group of automorphisms and $S\subseteq F$ a subfield. If for every $x\in F$ some nonzero $f\in\Z G$ satisfies $x^f\in S$, must $F=S$ or $F/S$ be purely inseparable?}
\expandafter\def\csname knscope@12.69\endcsname{\noindent\emph{Scope.} This is retained as a formulation issue, not as a substantive solution of an inferred intended question.}
\expandafter\def\csname knauthor@13.19\endcsname{Yu. M. Gorchakov}
\expandafter\def\csname knquestion@13.19\endcsname{Let $H\trianglelefteq Q\leq G_1\times\cdots\times G_n$, with both $H$ and $Q$ projecting onto every factor. If $Q/H$ is a $p$-group, must it be regular?}
\expandafter\def\csname knscope@13.19\endcsname{}
\expandafter\def\csname knauthor@13.42\endcsname{A. G. Myasnikov, V. N. Remeslennikov}
\expandafter\def\csname knquestion@13.42\endcsname{Show that the tensor $A$-completion of a free nilpotent group can fail to be nilpotent.}
\expandafter\def\csname knscope@13.42\endcsname{\noindent\emph{Scope.} Tensor completion is used in the unrestricted category specified by the cited definition.}
\expandafter\def\csname knauthor@14.22\endcsname{G. Baumslag, A. G. Myasnikov, V. N. Remeslennikov}
\expandafter\def\csname knquestion@14.22\endcsname{Must every irreducible system $S(X)=1$ over a torsion-free linear group $G$ be equivalent over $G$ to a finite system $T(X)=1$ with $\operatorname{Rad}_G(T)=\sqrt T$? Here $G[X]=G*F(X)$, the radical consists of words vanishing on every solution, and $\sqrt T$ is the least normal isolated subgroup containing $T$.}
\expandafter\def\csname knscope@14.22\endcsname{\noindent\emph{Scope.} The printed question does not require $G$ to be finitely generated.}
\expandafter\def\csname knauthor@14.26\endcsname{A. I. Budkin}
\expandafter\def\csname knquestion@14.26\endcsname{Is the quasivariety generated by all torsion-free nilpotent groups closed under direct wreath products $G\wr\Z$?}
\expandafter\def\csname knscope@14.26\endcsname{\noindent\emph{Scope.} Direct wreath product has the restricted-base meaning in the cited convention.}
\expandafter\def\csname knauthor@14.72\endcsname{K. N. Ponomar\"ev}
\expandafter\def\csname knquestion@14.72\endcsname{Let a finite cyclic group $G$ act on a regular algebraic variety $X$ over a field of arbitrary characteristic. If the fixed-point variety $X^G$ is a regular hypersurface, must $X/G$ be regular?}
\expandafter\def\csname knscope@14.72\endcsname{}
\expandafter\def\csname knauthor@15.65\endcsname{P. M. Neumann}
\expandafter\def\csname knquestion@15.65\endcsname{For the finite unitary, symplectic and orthogonal groups, are the limiting proportions, as dimension tends to infinity, of separable elements and of cyclic elements rational functions of the field size $q$? Separable means that the minimal polynomial has no repeated roots; cyclic means that the minimal and characteristic polynomials coincide.}
\expandafter\def\csname knscope@15.65\endcsname{}
\expandafter\def\csname knauthor@15.76\endcsname{B. I. Plotkin}
\expandafter\def\csname knquestion@15.76\endcsname{For a variety $\Theta$ of soluble groups, in particular metabelian groups, is every automorphism of the category $\Theta^0$ of finite-rank free groups inner, meaning naturally isomorphic to the identity?}
\expandafter\def\csname knscope@15.76\endcsname{\noindent\emph{Scope.} Part (b). The theorem below treats full varieties of fixed derived length, not all proper subvarieties; the distinction and a counterexample to the broader reading are retained.}
\expandafter\def\csname knauthor@15.89\endcsname{V. I. Trofimov}
\expandafter\def\csname knquestion@15.89\endcsname{For every infinite connected undirected vertex-transitive graph of finite valency without loops or multiple edges, is every complex number an eigenvalue of adjacency acting on all complex-valued functions on its vertices?}
\expandafter\def\csname knscope@15.89\endcsname{}
\expandafter\def\csname knauthor@15.92\endcsname{J. Wiegold}
\expandafter\def\csname knquestion@15.92\endcsname{For every $n>6$, does the triangle group with parameters $(2,3,n)$ have uncountably many pairwise nonisomorphic quotients that are residually finite alternating groups?}
\expandafter\def\csname knscope@15.92\endcsname{\noindent\emph{Scope.} The displayed Notebook presentation equates $x^2,y^3,(xy)^n$ without explicitly setting them to $1$. The proof uses the ordinary triangle-group quotient, so also covers that literal presentation.}
\expandafter\def\csname knauthor@16.9\endcsname{V. G. Bardakov, V. A. Tolstykh, V. E. Shpilrain}
\expandafter\def\csname knquestion@16.9\endcsname{Is there an algorithm computing the palindromic length of an element of a free group with a specified finite free basis? A palindrome is a reduced word read identically in either direction; palindromic length is the least number of palindromes whose product is the element.}
\expandafter\def\csname knscope@16.9\endcsname{\noindent\emph{Scope.} Qualitative computability is a prior consequence; the explicit reduction and operation bound below have separate attribution.}
\expandafter\def\csname knauthor@16.14\endcsname{Ya. G. Berkovich}
\expandafter\def\csname knquestion@16.14\endcsname{If a finite $2$-group $G$ satisfies $\Omega_1(G)\leq Z(G)$, is the rank of $G/G^2$ at most twice the rank of $Z(G)$?}
\expandafter\def\csname knscope@16.14\endcsname{}
\expandafter\def\csname knauthor@16.20\endcsname{A. I. Budkin}
\expandafter\def\csname knquestion@16.20\endcsname{For $H\leq A\in\mathcal M$, where $\mathcal M$ is a quasivariety, let $\operatorname{dom}_{\mathcal N}^{A}(H)$ consist of elements on which any two maps from $A$ to the same group in $\mathcal N$ agree whenever they agree on $H$. If the family of these dominions over all subquasivarieties $\mathcal N\subseteq\mathcal M$ is a lattice under inclusion, can it be modular and nondistributive?}
\expandafter\def\csname knscope@16.20\endcsname{\noindent\emph{Scope.} The convention concerning $A\in\mathcal N$ is discussed explicitly below.}
\expandafter\def\csname knauthor@16.28\endcsname{E. P. Vdovin}
\expandafter\def\csname knquestion@16.28\endcsname{Let $G$ be a connected reductive linear algebraic group over a field of positive characteristic and $X$ a closed subset. For $X^k=\{x_1\cdots x_k:x_i\in X\}$, must there exist an integer $c>1$ such that $X^c$ is closed?}
\expandafter\def\csname knscope@16.28\endcsname{\noindent\emph{Scope.} Part (a) only. The question allows a reducible closed set; part (b), concerning a conjugacy class, is not addressed.}
\expandafter\def\csname knauthor@16.38\endcsname{F. de Giovanni}
\expandafter\def\csname knquestion@16.38\endcsname{If $A,B$ are periodic subgroups of a soluble group $G$, must every subgroup contained in the set $AB$ be periodic?}
\expandafter\def\csname knscope@16.38\endcsname{}
\expandafter\def\csname knauthor@16.60\endcsname{D. MacHale}
\expandafter\def\csname knquestion@16.60\endcsname{For a finite group $G$, put $T(G)=\sum_{\chi\in\operatorname{Irr}(G)}\chi(1)$ and $S_\alpha=\{g:\alpha(g)=g^{-1}\}$ for $\alpha\in\Aut G$. Is $T(G)\geq|S_\alpha|$ always?}
\expandafter\def\csname knscope@16.60\endcsname{}
\expandafter\def\csname knauthor@16.87\endcsname{E. I. Timoshenko}
\expandafter\def\csname knquestion@16.87\endcsname{Let $G_r$ be the relatively free group of rank $r$ in a variety $\mathcal M$. A test set is a subset fixed pointwise only by automorphisms among endomorphisms; its least possible size is the test rank. If this rank equals $r$ for every $r\geq1$, (a) must $\mathcal M$ be abelian, and (b) if $\mathcal M$ is nonperiodic, must it be the variety of all abelian groups?}
\expandafter\def\csname knscope@16.87\endcsname{}
\expandafter\def\csname knauthor@17.33\endcsname{A. I. Budkin}
\expandafter\def\csname knquestion@17.33\endcsname{Can the quasivariety generated by $\langle a,b\mid a^{-1}ba=b^{-1}\rangle$ be axiomatized by quasi-identities using a fixed finite set of variables?}
\expandafter\def\csname knscope@17.33\endcsname{}
\expandafter\def\csname knauthor@17.34\endcsname{A. I. Budkin}
\expandafter\def\csname knquestion@17.34\endcsname{In the quasivariety of torsion-free nilpotent groups of class at most $c$, is every divisible subgroup equal to its dominion in each containing group of that quasivariety?}
\expandafter\def\csname knscope@17.34\endcsname{}
\expandafter\def\csname knauthor@17.39\endcsname{A. F. Vasil'ev, L. A. Shemetkov}
\expandafter\def\csname knquestion@17.39\endcsname{Is there a fixed positive integer $n$ such that the hypercenter of every finite soluble group is the intersection of $n$ system normalizers? If so, what is the least such $n$?}
\expandafter\def\csname knscope@17.39\endcsname{}
\expandafter\def\csname knauthor@17.101\endcsname{B. I. Plotkin}
\expandafter\def\csname knquestion@17.101\endcsname{Is every representation of a group embeddable into a homogeneous representation? Homogeneity means that every isomorphism between finitely generated subrepresentations extends to an automorphism. A representation $(V,G)$ is finitely generated when $G$ is finitely generated and $V$ is finitely generated over its group algebra.}
\expandafter\def\csname knscope@17.101\endcsname{}
\expandafter\def\csname knauthor@17.118\endcsname{E. I. Khukhro}
\expandafter\def\csname knquestion@17.118\endcsname{If a finite $p$-group has a subgroup of exponent $p$ and index $p$, must it have a characteristic subgroup of exponent $p$ and index bounded in terms of $p$ alone?}
\expandafter\def\csname knscope@17.118\endcsname{\noindent\emph{Scope.} Bounds for $p=2,3$ and bounds depending on nilpotency class are proved below; the unrestricted bound for $p\geq5$ is not established.}
\expandafter\def\csname knauthor@18.18\endcsname{O. V. Belegradek}
\expandafter\def\csname knquestion@18.18\endcsname{Are both the set of first-order sentences true in almost all finite groups and its complement non-computably-enumerable?}
\expandafter\def\csname knscope@18.18\endcsname{\noindent\emph{Scope.} Almost all means all but finitely many isomorphism types.}
\expandafter\def\csname knauthor@18.76\endcsname{E. A. Palyutin}
\expandafter\def\csname knquestion@18.76\endcsname{Let $A$ be a division ring and $G\leq A^\times$. Must every extension of the additive group of $A$ by $G$, inducing the multiplication action of $G$ on $A$, split?}
\expandafter\def\csname knscope@18.76\endcsname{}
\expandafter\def\csname knauthor@18.92\endcsname{A. N. Skiba}
\expandafter\def\csname knquestion@18.92\endcsname{A complete lattice of formations means a nonempty family of formations closed under arbitrary intersections and having a largest member. Do such lattices of formations of finite groups exist which are (a) nonalgebraic or (b) nonmodular? Algebraic means that every element is a join of compact elements, and compactness refers to finite subjoins.}
\expandafter\def\csname knscope@18.92\endcsname{\noindent\emph{Scope.} Joins are taken within the specified family of formations.}
\expandafter\def\csname knauthor@18.120\endcsname{E. Jabara}
\expandafter\def\csname knquestion@18.120\endcsname{Let a finite $p$-group be $P=AB$, with $A$ abelian and $B$ of class two; assume $A\cap B=1$ if necessary. Must $\langle A,B^\prime\rangle=AB_0$ for an abelian subgroup $B_0\leq B$?}
\expandafter\def\csname knscope@18.120\endcsname{\noindent\emph{Scope.} The argument below proves the case of total nilpotency class at most four, without requiring finiteness or trivial intersection.}
\expandafter\def\csname knauthor@19.9\endcsname{V. G. Bardakov, M. V. Neshchadim}
\expandafter\def\csname knquestion@19.9\endcsname{For a finitely generated linear group $G$, is its nonabelian tensor square $G\otimes G$ linear? In particular, is this true for the braid groups $B_n$, $n>3$?}
\expandafter\def\csname knscope@19.9\endcsname{\noindent\emph{Scope.} The scope of each prior theorem and construction is stated below; no general preservation under finite central extensions is assumed.}
\expandafter\def\csname knauthor@19.48\endcsname{V. A. Koibaev}
\expandafter\def\csname knquestion@19.48\endcsname{Let $R$ be a principal ideal domain, $K$ an algebraic extension of its fraction field, and $\sigma$ an irreducible elementary net over $K$ whose levels are $R$-modules. Must $\sigma$ be closed? Net levels satisfy $\sigma_{ir}\sigma_{rj}\subseteq\sigma_{ij}$; irreducible means every off-diagonal level is nonzero, and closed means its elementary net subgroup contains no additional transvections.}
\expandafter\def\csname knscope@19.48\endcsname{}
\expandafter\def\csname knauthor@19.56\endcsname{V. S. Monakhov}
\expandafter\def\csname knquestion@19.56\endcsname{A $\star$-commutator is $[x,y]$ with $x,y$ of coprime prime-power orders. Is a finite group $G$ soluble whenever $|ab|\geq|a||b|$ for any two $\star$-commutators $a,b$ of coprime orders?}
\expandafter\def\csname knscope@19.56\endcsname{}
\expandafter\def\csname knauthor@19.61\endcsname{Ya. N. Nuzhin}
\expandafter\def\csname knquestion@19.61\endcsname{Let $A=(A_r)$ be an elementary carpet over a commutative ring $K$, with carpet subgroup $\Phi(A)=\langle x_r(A_r)\rangle$. Set $\overline A_r=\{t\in K:x_r(t)\in\Phi(A)\}$. Is $\overline A=(\overline A_r)$ always a carpet?}
\expandafter\def\csname knscope@19.61\endcsname{}
\expandafter\def\csname knauthor@19.62\endcsname{Ya. N. Nuzhin}
\expandafter\def\csname knquestion@19.62\endcsname{For an elementary carpet $A=(A_r)_{r\in\Phi}$ of rank at least two over a commutative ring, form its derived carpet $B_p=\sum_{ir+js=p}C_{ij,rs}A_r^iA_s^j$ using the Chevalley commutator constants. Is every such derived carpet closed, that is, does its carpet subgroup contain no new root elements?}
\expandafter\def\csname knscope@19.62\endcsname{}
\expandafter\def\csname knauthor@19.63\endcsname{Ya. N. Nuzhin}
\expandafter\def\csname knquestion@19.63\endcsname{For an elementary carpet $A=(A_r)$ over a commutative ring, do the inclusions $A_r^2A_{-r}\subseteq A_r$ for every root, where $A_r^2=\{t^2:t\in A_r\}$, imply that the carpet is closed?}
\expandafter\def\csname knscope@19.63\endcsname{\noindent\emph{Scope.} The supplied Notebook already marks this criterion solved by Nuzhin. Its use below carries prior-work status and adds no candidate.}
\expandafter\def\csname knauthor@19.83\endcsname{P. Shumyatsky}
\expandafter\def\csname knquestion@19.83\endcsname{Is the set of almost Engel elements of a linear group a subgroup? An element $g$ is almost Engel if it has a finite set $E(g)$ such that, for each $x$, all sufficiently long left-normed commutators $[x,g,\ldots,g]$ belong to $E(g)$.}
\expandafter\def\csname knscope@19.83\endcsname{}
\expandafter\def\csname knauthor@19.93\endcsname{P. Spiga}
\expandafter\def\csname knquestion@19.93\endcsname{Is there a function $f(p,d)$ bounding the order of every factor of the lower $p$-central series of every $d$-generated finite $p$-group with no quotient isomorphic to $C_p\wr C_p$?}
\expandafter\def\csname knscope@19.93\endcsname{}
\expandafter\def\csname knauthor@19.108\endcsname{T. Wilde}
\expandafter\def\csname knquestion@19.108\endcsname{For a finite $p$-group $P$, with $p$ odd, and a complex irreducible character $\chi$, does $\chi(x)\ne0$ imply $|x|\mid |P|/\chi(1)^2$?}
\expandafter\def\csname knscope@19.108\endcsname{\noindent\emph{Scope.} The additional hypothesis on a large abelian normal subgroup in the order-$p^{10}$ case is essential to the result below.}
\expandafter\def\csname knauthor@20.1\endcsname{M. Arezoomand, M. A. Iranmanesh, C. E. Praeger, G. Tracey}
\expandafter\def\csname knquestion@20.1\endcsname{Does there exist a finite insoluble totally $2$-closed group with nontrivial Fitting subgroup? Totally $2$-closed means that in every faithful permutation action the group is the full group preserving each of its orbits on ordered pairs.}
\expandafter\def\csname knscope@20.1\endcsname{}
\expandafter\def\csname knauthor@20.8\endcsname{G. M. Bergman}
\expandafter\def\csname knquestion@20.8\endcsname{Let $K<H<F$ be finite-rank free groups, with $\operatorname{rank}H<\operatorname{rank}K$, but every proper subgroup of $H$ containing $K$ of rank at least $\operatorname{rank}K$. Is inclusion the only homomorphism $H\to F$ fixing $K$ pointwise?}
\expandafter\def\csname knscope@20.8\endcsname{}
\expandafter\def\csname knauthor@20.11\endcsname{G. M. Bergman}
\expandafter\def\csname knquestion@20.11\endcsname{Let $F\leq H$ be free groups with a finite-rank intermediate subgroup, and let $r$ be the least possible intermediate rank. Must there be (i) a largest or (ii) a smallest intermediate group of rank $r$? Must respectively (i\textquotesingle) the join or (ii\textquotesingle) the intersection of any two rank-$r$ intermediate groups have rank $r$?}
\expandafter\def\csname knscope@20.11\endcsname{}
\expandafter\def\csname knauthor@20.33\endcsname{A. Bauer; communicated by J. Grochow}
\expandafter\def\csname knquestion@20.33\endcsname{For every oracle $X\subseteq\mathbb N$, does there exist a finitely generated $X$-computably-enumerably presented group $G_X$ such that a finitely generated group has such a presentation exactly when it embeds in a quotient of a finite free product of copies of $G_X$ by the normal closure of finitely many elements?}
\expandafter\def\csname knscope@20.33\endcsname{}
\expandafter\def\csname knauthor@20.89\endcsname{P. Shumyatsky}
\expandafter\def\csname knquestion@20.89\endcsname{Is the set of almost Engel elements of a linear group a subgroup? An element $g$ is almost Engel if it has a finite set $E(g)$ such that, for each $x$, all sufficiently long left-normed commutators $[x,g,\ldots,g]$ belong to $E(g)$.}
\expandafter\def\csname knscope@20.89\endcsname{\noindent\emph{Scope.} Also Problem 19.83, by the same author. General positive characteristic remains unresolved here.}
\expandafter\def\csname knauthor@20.90\endcsname{P. Shumyatsky}
\expandafter\def\csname knquestion@20.90\endcsname{Find an infinite finitely generated profinite group, not prosoluble, in which the centralizer of every nonidentity element is pronilpotent.}
\expandafter\def\csname knscope@20.90\endcsname{}
\expandafter\def\csname knauthor@20.92\endcsname{A. Smoktunowicz}
\expandafter\def\csname knquestion@20.92\endcsname{Let $A$ be an $\F_p$-brace of order $p^k$, and write $a*b=a\circ b-a-b$, with $a*(\alpha b)=\alpha(a*b)$. For $p$ sufficiently large relative to $k$, does $a\cdot b=-\sum_{i=0}^{p-2}2^{-i}((2^i a)*b)$ make $A$ a pre-Lie algebra?}
\expandafter\def\csname knscope@20.92\endcsname{\noindent\emph{Scope.} Part (a) only. The negative sign, full range of summation and absence of a strong-nilpotence hypothesis are essential. Part (b) is already marked solved in the supplied Notebook.}
\expandafter\def\csname knauthor@20.100\endcsname{Z.-W. Sun}
\expandafter\def\csname knquestion@20.100\endcsname{If $G$ has no elements of orders $2,\ldots,n+1$, must every $n$-element subset admit an ordering $a_1,\ldots,a_n$ with $a_1,a_2^2,\ldots,a_n^n$ pairwise distinct?}
\expandafter\def\csname knscope@20.100\endcsname{\noindent\emph{Scope.} The unrestricted question remains open here; the reductions, torsion-free case and certified small values of $n$ are stated below.}
\expandafter\def\csname knauthor@20.108\endcsname{C. Tsang}
\expandafter\def\csname knquestion@20.108\endcsname{Put $T(G)=\NHol(G)/\Hol(G)$, where $\Hol(G)$ is the normalizer of the left translations in $\operatorname{Sym}(G)$ and $\NHol(G)$ its normalizer. Is there a centerless group $G$, and in particular a finite one, for which $T(G)$ is not an elementary abelian $2$-group?}
\expandafter\def\csname knscope@20.108\endcsname{\noindent\emph{Scope.} Parts (a,b). Part (c), attributed to A. Caranti, is a separate $p$-group question and is not answered here.}
\expandafter\def\csname knauthor@20.122\endcsname{V. I. Zenkov}
\expandafter\def\csname knquestion@20.122\endcsname{For nilpotent subgroups $A,B,C$ of a finite group $G$, let $\operatorname{Min}_G(A,B,C)$ be generated by inclusion-minimal intersections $A\cap B^x\cap C^y$, and let $\operatorname{min}_G(A,B,C)$ be generated by those of minimum order. Must (a) $\operatorname{min}_G(A,B,C)$ or (b) $\operatorname{Min}_G(A,B,C)$ lie in $F(G)$? (c) Do these assertions hold when $G$ is soluble?}
\expandafter\def\csname knscope@20.122\endcsname{}
\expandafter\def\csname knauthor@20.124\endcsname{V. G. Bardakov}
\expandafter\def\csname knquestion@20.124\endcsname{Does every nonabelian free group $F$ admit a Rota--Baxter operator $B:F\to F$ with image $F^\prime$? Such an operator satisfies $B(g)B(h)=B(gB(g)hB(g)^{-1})$.}
\expandafter\def\csname knscope@20.124\endcsname{\noindent\emph{Scope.} The construction below covers infinite rank; finite rank remains unresolved here.}
\expandafter\def\csname knauthor@21.40\endcsname{A. Dantas, E. de Melo}
\expandafter\def\csname knquestion@21.40\endcsname{Is every subgroup of $\GL(n,\Q)$ having finitely many automorphism orbits virtually soluble?}
\expandafter\def\csname knscope@21.40\endcsname{}
\expandafter\def\csname knauthor@21.42\endcsname{A. Dantas, S. Sidki}
\expandafter\def\csname knquestion@21.42\endcsname{Is there a three-generated torsion-free nilpotent group of class three which is not self-similar?}
\expandafter\def\csname knscope@21.42\endcsname{}
\expandafter\def\csname knauthor@21.60\endcsname{D. Johnston, D. Rumynin}
\expandafter\def\csname knquestion@21.60\endcsname{For a finite group $G$, let $\mathcal X$ consist of the $\F_pG$-modules obtained by reducing simple $\Q G$-modules. Is $\Z_{(p)}G$ semiperfect if and only if each projective indecomposable $\F_pG$-module is an $\mathbb N$-linear combination of members of $\mathcal X$ in the Grothendieck group?}
\expandafter\def\csname knscope@21.60\endcsname{}
\expandafter\def\csname knauthor@21.68\endcsname{M. Kida}
\expandafter\def\csname knquestion@21.68\endcsname{Is every finite semi-abelian group monomial? Semi-abelian means that there is a chain $1=G_0\leq\cdots\leq G_n=G$ such that each $G_{i+1}$ is isomorphic to a quotient of $A_i\rtimes G_i$ for an abelian group $A_i$.}
\expandafter\def\csname knscope@21.68\endcsname{}
\expandafter\def\csname knauthor@21.76\endcsname{V. A. Koibaev}
\expandafter\def\csname knquestion@21.76\endcsname{Over fields of odd characteristic, do there exist irreducible closed elementary nets of order $n\geq3$ which cannot be completed by diagonal entries? Irreducible means all off-diagonal levels are nonzero; closed means the elementary net group contains no new transvections.}
\expandafter\def\csname knscope@21.76\endcsname{}
\expandafter\def\csname knauthor@21.106\endcsname{M. Petschick}
\expandafter\def\csname knquestion@21.106\endcsname{Is every first-order group formula with one free variable concise in residually finite groups? Here concise means that whenever its truth set in a group is finite, the subgroup generated by that set is finite.}
\expandafter\def\csname knscope@21.106\endcsname{}
\expandafter\def\csname knauthor@21.107\endcsname{I. V. Protasov}
\expandafter\def\csname knquestion@21.107\endcsname{If a countable topological group can be partitioned into countably many dense subsets, must it have an expansive sequence? Such a sequence consists of pairwise disjoint finite sets $F_n$ meeting every nonempty open set for all sufficiently large $n$.}
\expandafter\def\csname knscope@21.107\endcsname{}
\expandafter\def\csname knauthor@21.114\endcsname{L. Sabatini}
\expandafter\def\csname knquestion@21.114\endcsname{Is the derived length of finite weakly ab-maximal groups bounded? Weakly ab-maximal means $|H:H^\prime|\leq|G:G^\prime|$ for every $H\leq G$.}
\expandafter\def\csname knscope@21.114\endcsname{\noindent\emph{Scope.} The construction below gives derived length three, not an unbounded family.}
\expandafter\def\csname knauthor@21.115\endcsname{B. Sambale}
\expandafter\def\csname knquestion@21.115\endcsname{If $U=C_1\cup\cdots\cup C_n\ne G$ is a union of left or right cosets in a finite group $G$, must $|G\setminus U|\geq |G|/2^n$?}
\expandafter\def\csname knscope@21.115\endcsname{}
\expandafter\def\csname knauthor@21.121\endcsname{C. Shramov}
\expandafter\def\csname knquestion@21.121\endcsname{A group $\Gamma$ is $p$-Jordan if every finite $G\leq\Gamma$ has a normal abelian subgroup of order coprime to $p$ and index at most $J|G_{(p)}|^e$, for some constants $J,e$. Let $e(\Gamma)$ be the infimum of admissible $e$. (a) Is the infimum attained? (b) Is $e(\Gamma)\leq3$ always?}
\expandafter\def\csname knscope@21.121\endcsname{\noindent\emph{Scope.} Part (a) is answered in Section~\ref{cand:21.121}; only partial results for (b) are claimed in Section~\ref{partial:jordan}.}
\expandafter\def\csname knauthor@21.132\endcsname{A. V. Timofeenko}
\expandafter\def\csname knquestion@21.132\endcsname{For every prime $p$, construct, using the development of Golod\textquotesingle s construction, a finitely generated residually finite $p$-group with nontrivial finite center.}
\expandafter\def\csname knscope@21.132\endcsname{}
